% Generated draft: requires mathematical and source review.
\documentclass[11pt]{article}
\usepackage[margin=25mm]{geometry}
\usepackage{fontspec}
\setmainfont{Latin Modern Roman}
\usepackage{amsmath,amssymb,mathtools,unicode-math}
\setmathfont{Latin Modern Math}
\usepackage{enumitem,needspace,xurl,hyperref}
\setlength{\parindent}{0pt}
\setlength{\parskip}{5pt}
\setlength{\emergencystretch}{3em}
\newcommand{\N}{\mathbb{N}}
\newcommand{\Z}{\mathbb{Z}}
\newcommand{\Q}{\mathbb{Q}}
\newcommand{\R}{\mathbb{R}}
\newcommand{\C}{\mathbb{C}}
\newcommand{\F}{\mathbb{F}}
\newcommand{\A}{\mathbb{A}}
\newcommand{\PP}{\mathbb{P}}
\newcommand{\E}{\mathbb{E}}
\newcommand{\eps}{\varepsilon}
\newcommand{\dd}{\,\mathrm{d}}
\begin{document}
\section*{UnsolvedMath ID: 11000071}
\textbf{Source problem code:} catalogue entry
\subsection{Definitions and mathematical statement}
Let \(\operatorname{Mod}_{g,1}\) denote the source's mapping class group and let \(\rho_{\mathrm{std}}:\operatorname{Mod}_{g,1}\to\operatorname{Homeo}^+(S^1)\) be the standard action referenced in its missing equation (19). The assertion that this group has torsion indicates a marked-point or puncture convention rather than a boundary-fixed convention, but the exact definition and genus range must be recovered.

\emph{Question 6.2.} Is every faithful action \(\rho:\operatorname{Mod}_{g,1}\to\operatorname{Homeo}^+(S^1)\) conjugate to \(\rho_{\mathrm{std}}\) by an orientation-preserving circle homeomorphism? Explicitly, must there be \(h\in\operatorname{Homeo}^+(S^1)\) with
\[
 \rho(\gamma)=h\rho_{\mathrm{std}}(\gamma)h^{-1}
 \quad\text{for every }\gamma\in\operatorname{Mod}_{g,1}?
\]
Ask the same question for every finite-index subgroup \(\Gamma\leq\operatorname{Mod}_{g,1}\), comparing a faithful action of \(\Gamma\) with \(\rho_{\mathrm{std}}|_\Gamma\). The source asks exact conjugacy in \(\operatorname{Homeo}^+(S^1)\), not semiconjugacy or conjugacy allowing reversal of the circle orientation, and does not impose minimality. These distinctions are essential and are not silently weakened.

The source proposes a stronger dynamical characterization analogous to the Convergence Groups Theorem of Tukia, Casson--Jungreis, and Gabai, replacing asymptotic source--sink behavior by alternating source--sink--\(\cdots\)--source--sink behavior. No precise sequential condition, number of alternations, or refinement is specified, so this remains a suggested research direction.

Its context describes the lift group
\[
 \widetilde{\operatorname{Homeo}}^+(S^1)
 =\{f\in\operatorname{Homeo}^+(\R):f(x+1)=f(x)+1\},
\]
with the central extension by integer translations. Pullback through \(\rho_{\mathrm{std}}\) gives
\[
 1\longrightarrow\Z\longrightarrow
 \widetilde{\operatorname{Mod}}_{g,1}
 \longrightarrow\operatorname{Mod}_{g,1}\longrightarrow1.
 \tag{20}
\]
The source uses the torsion-free lift group and torsion in \(\operatorname{Mod}_{g,1}\) to say (20) does not split, giving a nonzero class \(\xi\in H^2(\operatorname{Mod}_{g,1};\Z)\), identified with the Euler cocycle. It describes evaluation using pointed maps \(\Sigma_h\to\mathcal M_g\) and a pulled-back ``universal circle bundle.'' The relevant orbifold, bundle, and pointed/unpointed base conventions are not supplied, so that geometric interpretation requires review.

Finally, it notes that torsion obstructs left-orderings and faithful orientation-preserving actions on \(\R\) for these mapping class groups, says this is the only obstruction then known to the author, and notes that torsion disappears in appropriate finite-index subgroups. The last statement does not assert that the resulting torsion-free subgroups are left-orderable or have faithful line actions.
\subsection{Short English statement}
Ask whether every faithful circle action of the pointed mapping class group, and of its finite-index subgroups, is conjugate by an orientation-preserving homeomorphism to the specified standard action. Exact conjugacy, semiconjugacy, and minimal-action variants must be distinguished; the source also records the lift extension and Euler-class context.
\subsection{Sources}
\begin{itemize}
\item Local UnsolvedMath registry, permanent ID 11000071, source code catalogue entry, Question 6.2, ``Rigidity of the \(\operatorname{Mod}_{g,1}\) action on \(S^1\),'' PDF page 54. No individual proposer is supplied.
\item Farb, \emph{Problems on Mapping Class Groups and Related Topics}. Supplied URL: \url{https://www.math.uchicago.edu/~farb/papers/mcgbook.pdf}; recorded access 2026-07-29; extraction pdf-font-structure. The volume contains 329 explicitly typeset headings across chapters, with repeated numbers disambiguated by page. The defining equation (19) is absent; the central extension is labeled (20).
\item The statement credits Tukia, Casson--Jungreis, and Gabai for the convergence-groups theorem. The summary mentions Matsumoto and unspecified circle-action results, with an unresolved genus placeholder, but gives no complete theorem or bibliography.
\item Supplied status open and OPEN-TRIAGE. Its summary says the specific action-rigidity formulation was not fully verified. Historical phrases such as ``as far as we know'' about line actions are not treated as a 2026 audit. Current primary-literature status and redistribution terms remain NEEDS\_REVIEW.
\end{itemize}
\end{document}
